\documentclass[10pt,a4paper]{amsart}
\usepackage[margin=19mm]{geometry}
\usepackage{amsmath,amssymb,mathtools}
\usepackage[scr=rsfs,cal=euler]{mathalfa}
\usepackage{xfrac}
\usepackage[unicode,hidelinks,bookmarks=false]{hyperref}
\newcommand{\ie}{i.e.\ }
\newcommand{\cf}{cf.\ }
\newcommand{\resp}{respectively\ }
\newcommand{\Subsection}{\subsection}
\providecommand{\nobreakdash}{\nobreak\hbox{-}}
\setlength{\emergencystretch}{2em}
\multlinegap0pt
\usepackage{amssymb}
\usepackage{mathtools}
\usepackage[shortlabels]{enumitem}
\usepackage{float}
\newtheorem{proposition}{Proposition}
\title{Combinatorial aspects of normal ordering of \\ 3-dimensional skew polynomial rings}
\author{Andr\'es Rubiano, Armando Reyes}
\date{Source: arXiv:2605.03094v2 (CC BY 4.0)}
\begin{document}
\maketitle

\noindent\emph{Source and licence.} Combinatorial aspects of normal ordering of \\ 3-dimensional skew polynomial rings. Version arXiv:2605.03094v2 (CC BY 4.0), distributed by arXiv under the Creative Commons Attribution 4.0 International licence, http://creativecommons.org/licenses/by/4.0/, as recorded in the arXiv licence metadata for this exact version and shown on the abstract page https://arxiv.org/abs/2605.03094v2. Full licence notice: \texttt{licenses/arxiv-2605.03094-CC-BY-4.0.txt}.\par\smallskip

\noindent\textit{Editorial scope.} Counted unit: the article's proposition case(2)(ii) (source0.tex lines 635--719). The article's complete original proof of it is delivered with it (source0.tex lines 720--807, one \texttt{proof} environment of 88 lines). No mathematics is written, repaired or re-derived: the statement and the proof are the article's own lines, in the article's own notation, and no retained line is substituted or re-flowed.  No further source-local context is needed: every cross-reference of every retained line resolves inside the two delivered spans.  Nothing was omitted from any retained span, so the package carries no omission record.  No retained line contains a citation, so the wrapper carries no bibliography and no bibliography entry.  No retained line contains a figure environment, an image-inclusion command or drawing code, so no image is delivered and the package declares no asset dependency.  Editorial formatting only, in addition: an AMS \texttt{amsart} preamble that restates the article's own declarations and macros used by the retained lines, namely the environments \texttt{proposition} on separate counters, together with the standard packages those lines need.  The retained environments print in this document as Proposition 1 in order of appearance, whereas the article prints the same items with the numbering of its own full document.  The complete native source of the article as distributed in the arXiv e-print for this exact version is bundled unchanged as \texttt{sources/199770/source0.tex}.
\begin{proposition}\label{case(2)(ii)}
Let $A$ be the $\Bbbk$-algebra generated by $x,y,z$ with relations
\[
yz - zy = z, \quad zx - \beta xz = b \quad {\rm and} \quad xy - yx = x,
\]
where $\beta\in\Bbbk^{\times}$ and $b\in\Bbbk$. Then:
\begin{enumerate}
\item [\rm (1)] For all $m,n\in\mathbb{N}$,
\[
y^n x^m = x^m (y-m)^n \quad {\rm and}\quad z^n y^m = (y-n)^m z^n.
\]
Furthermore, for all $m,n\in\mathbb{N}$ we have the closed PBW normal form given by
\begin{equation}\label{znxmclosedconstb2ii}
z^n x^m = \sum_{k=0}^{\min\{m,n\}}
\beta^{(m-k)(n-k)}b^k\begin{bmatrix} n\\ k\end{bmatrix}_{\beta}\begin{bmatrix} m\\ k\end{bmatrix}_{\beta}[k]_{\beta}!x^{m-k}z^{n-k}.
\end{equation}
Define the elements $W^{(m)}_{n,k}\in\Bbbk$ as follows: 
\begin{equation}\label{Wdefconstb2ii}
z^n x^m = \sum_{k=0}^{\min\{m,n\}} x^{m-k}W^{(m)}_{n,k}z^{n-k}.
\end{equation}
Then $W^{(m)}_{0,0}=1$, $W^{(m)}_{0,k}=0$ $(k\ge 1)$ and we get the recurrence relation given by
\begin{equation}\label{Wrecconstb2ii}
W^{(m)}_{n+1,k}=\beta^{m-k}W^{(m)}_{n,k}+b[m-k+1]_{\beta}W^{(m)}_{n,k-1}, \quad  \text{for } n\ge 0,\ k\ge 0,
\end{equation}
with $W^{(m)}_{n,-1}:=0$.

\item [\rm (2)] For all $m, n, s \in \mathbb{N}$, we have that
\[
(x^n y^m)^s = x^{ns}\prod_{j=0}^{s-1}(y-jn)^m \quad {\rm and} \quad (y^n z^m)^s = \left(\prod_{j=0}^{s-1}(y-jm)^n\right)z^{ms}.
\]
Moreover, $(x^n z^m)^s$ has a unique PBW expansion as 
\begin{equation}\label{xnzmUconstb2ii}
(x^n z^m)^s = \sum_{\ell=0}^{\min\{ns,ms\}} x^{ns -\ell}U^{(n,m)}_{s,\ell}z^{ms-\ell}, \quad U^{(n,m)}_{s,\ell}\in\Bbbk,
\end{equation}
with initial conditions $U^{(n,m)}_{1,0}=1$, $U^{(n,m)}_{1,\ell}=0$ for $\ell\ge 1$ and the explicit recursion
\begin{equation}\label{Urecconstb2ii}
U^{(n,m)}_{s+1,\ell} = \sum_{k=0}^{\min\{n,\ell\}}U^{(n,m)}_{s,\ell-k}W^{(n)}_{ms-(\ell-k),k}, \quad \text{for } s\ge 1, \ell\ge 0,
\end{equation}
where $W^{(n)}_{r,k}$ are the normal ordering coefficients from \eqref{Wdefconstb2ii} {\rm (}with $m$ replaced by $n${\rm )}.

\item [\rm (3)] There exist uniquely determined polynomials $V_{s,\ell}(y)\in\Bbbk[y]$ for $0\le \ell\le s$ such that
\begin{equation}\label{Vdefconstb2ii}
(xyz)^s=\sum_{\ell=0}^{s} x^{s-\ell}V_{s,\ell}(y)z^{s-\ell}.
\end{equation}
All of them satisfy that $V_{1,0}(y) = y$, $V_{1,\ell}(y) = 0$ for $\ell\ge 1$ and the explicit recursion
\begin{equation}\label{Vrecconstb2ii}
V_{s+1,\ell}(y) =(y-(s-\ell))\Big(\beta^{s-\ell}V_{s,\ell}(y-1)+b[s-\ell+1]_\beta  V_{s,\ell-1}(y)\Big), 
\end{equation}
for $s\ge 1$ and $\ell\ge 0$, where $V_{s,-1}(y):=0$.

\item[\rm (4)] For all $n, m, t, s\in\mathbb{N}$ and  $0\le \ell\le \min\{ns,ts\}$, there exist uniquely determined polynomials
$R^{(n,m,t)}_{s,\ell}(y)\in\Bbbk[y]$,  such that
\begin{equation}\label{Rdefconstb2ii}
(x^n y^m z^t)^s = \sum_{\ell=0}^{\min\{ns,ts\}} x^{ns-\ell}R^{(n,m,t)}_{s,\ell}(y)z^{ts-\ell}.
\end{equation}
They satisfy $R^{(n,m,t)}_{1,0}(y)=y^m$, $R^{(n,m,t)}_{1,\ell}(y)=0$ for $\ell\ge 1$ and 
\begin{equation}\label{Rrecconstb2ii}
R^{(n,m,t)}_{s+1,\ell}(y)=(y-(ts-\ell))^{m} \sum_{k=0}^{\min\{n,\ell\}}R^{(n,m,t)}_{s,\ell-k}\big(y-(n-k)\big)
W^{(n)}_{ts-(\ell-k),k},
\end{equation}
for $s\ge 1$ and $\ell\ge 0$, where $W^{(n)}_{r,k}$ are as in \eqref{Wdefconstb2ii}.

\item[\rm (5)] There exist uniquely determined polynomials $S_{n,k}(y),T_{n,k}(y)\in\Bbbk[y]$ such that
\[
(x+y)^n = \sum_{k=0}^{n} x^{k}S_{n,k}(y) \quad {\rm and} \quad (y+z)^n=\sum_{k=0}^{n} T_{n,k}(y)z^{k},
\]
characterized by the following explicit recursions:
\[
S_{0,0}(y) = 1,\ S_{0,k}(y)=0\ (k\ge 1),\quad S_{n+1,k}(y)=S_{n,k-1}(y)+(y-k)S_{n,k}(y),
\]
and 
\[
T_{0,0}(y)=1,\ T_{0,k}(y)=0\ (k\ge 1),\quad T_{n+1,k}(y)=yT_{n,k}(y)+T_{n,k-1}(y-1).
\]
The expression $(x+z)^n$ has a unique PBW expansion given by
\begin{equation}\label{Edefconstb2ii}
(x+z)^n = \sum_{\substack{i,k\ge 0\\ i+k\le n}} x^{i}E_{n;i,k}z^{k}, \quad E_{n;i,k}\in\Bbbk[b],
\end{equation}
with $E_{0;0,0}=1$ and $E_{0;i,k}=0$ for $(i,k)\neq(0,0)$, and the explicit recursion
\begin{equation}\label{Erecconstb2ii}
E_{n+1;i,k}=E_{n;i,k-1}+\beta^{k}E_{n;i-1,k}+b[k+1]_\beta E_{n;i,k+1},
\end{equation}
where we set $E_{n;i,k} = 0$ whenever $i < 0$ or $k < 0$ or $i + k > n$.
\end{enumerate}
\end{proposition}

\begin{proof}
\begin{enumerate}
    \item [\rm (1)]  From $yx = xy-x$ we get $yx = x(y-1)$. For $f(y)\in\Bbbk[y]$,
\begin{equation}\label{shiftxcase2ii}
f(y)  x^m=x^m f(y-m) \text{ for } m\in\mathbb{N},
\end{equation}        
and taking $f(y) = y^n$ yields $y^n x^m=x^m(y-m)^n$.

From $zy = yz - z = (y-1)z$ we obtain that
\begin{equation}\label{shiftzcase2ii}
z^n f(y)=f(y-n)  z^n\text{ for } n\in\mathbb{N},\ f\in\Bbbk[y],
\end{equation}
and so $z^n y^m = (y - n)^m z^n$.

Next, using $zx = \beta xz+b$ and induction on $m$, it follows the identity
\begin{equation}\label{zxmcase2ii}
z x^m=\beta^m x^m z+b  [m]_\beta  x^{m-1}, \quad {\rm for} \, m\ge 1.
\end{equation}
Indeed, if \eqref{zxmcase2ii} holds for $m-1$, then
\begin{align*}
z x^m &\ =(zx)x^{m-1}=(\beta xz+b)x^{m-1} =\beta x(zx^{m-1})+b x^{m-1} \\
&\ =\beta^m x^m z+b(1+\beta+\cdots+\beta^{m-1})x^{m-1}.
\end{align*}
To prove \eqref{znxmclosedconstb2ii}, we proceed by induction on $n$ repeatedly using \eqref{zxmcase2ii}. Writing $z^{n+1}x^m = z(z^2n x^m)$ and expanding with \eqref{Wdefconstb2ii} yields a triangular recursion for the coefficients $W^{(m)}_{n,k}$, namely \eqref{Wrecconstb2ii}. Solving this recursion (or checking directly
by induction on $n$ using \eqref{zxmcase2ii} together with the $q$-Pascal identities for $[  \cdot  ]_\beta$)
gives the closed form \eqref{znxmclosedconstb2ii}. The alternative expression using Gaussian binomials follows from
\[
\frac{[r]^{\underline{k}}_\beta}{[k]_\beta!}
=\frac{[r]_\beta!}{[r-k]_\beta!  [k]_\beta!}
=\begin{bmatrix} r\\ k\end{bmatrix}_\beta.
\]

\item [\rm (2)] Using \eqref{shiftxcase2ii} with $f(y) = y^m$ yields $y^m x^n = x^n (y-n)^m$, and so
\begin{align*}
  (x^n y^m)^{s+1} = &\ (x^n y^m)^s x^n y^m  \\   
  = &\ x^{ns}\left(\prod_{j=0}^{s-1}(y-jn)^m\right)x^n y^m \\
  = &\ x^{n(s+1)}\prod_{j=0}^{s}(y-jn)^m, 
\end{align*}
which proves $(x^n y^m)^s$ by induction on $s$. The formula for $(y^n z^m)^s$ is analogous, using $z^m y^n = (y - m)^n z^m$ (which is \eqref{shiftzcase2ii}).

For $(x^n z^m)^s$, PBW gives the expansion \eqref{xnzmUconstb2ii}. Multiply \eqref{xnzmUconstb2ii} by $x^n z^m$ on the right and normal order each subword $z^{ms-(\ell-k)}x^n$ using \eqref{Wdefconstb2ii}. Since all coefficients are scalars (no $y$ occurs in $zx - \beta xz = b$), grouping the resulting terms by the total loss index yields the explicit recursion \eqref{Urecconstb2ii}.

\item [\rm (3)] Suppose that \eqref{Vdefconstb2ii} for some $s$. Multiply by $xyz$ on the right
\[
(xyz)^{s+1} = \sum_{\ell} x^{s-\ell}V_{s,\ell}(y)z^{s-\ell}xyz.
\]
and set $r:=s-\ell$. Using $z^{r}x=\beta^{r}x z^{r}+b[r]_\beta z^{r-1}$ (the case $m=1$ of \eqref{znxmclosedconstb2ii} or a direct induction from $zx=\beta xz+b$), we obtain two contributions.
For the first one, commute $V_{s,\ell}(y)$ past $x$ using $f(y)x=x f(y-1)$ and then commute $z^r$ past $y$ using
$z^r y=(y-r)z^r$ (from \eqref{shiftzcase2ii}), to get
\[
\beta^{r}(y-r)  x^{s-\ell+1}V_{s,\ell}(y-1)z^{r+1}.
\]
For the second one, use again $z^{r-1}y=(y-r+1)z^{r-1}$ and so 
\[
b[r]_\beta (y-r)x^{s-\ell}V_{s,\ell}(y)z^{r}.
\]
Reindexing the latter as a contribution to $V_{s+1,\ell}(y)$ coming from $V_{s,\ell-1}(y)$ yields exactly \eqref{Vrecconstb2ii}.

\item [\rm (4)] Write $B:=x^n y^m z^t$. Assuming \eqref{Rdefconstb2ii} for $s$, multiply by $B$ on the right:
\[
B^{s+1}=\sum_{\ell} x^{ns-\ell}R_{s,\ell}(y)  z^{ts-\ell}x^n y^m z^t 
\quad {\rm and} \quad R_{s,\ell}(y):=R^{(n,m,t)}_{s,\ell}(y).
\]
Put $r := ts - l$ and normal order $z^{r}x^n$ using \eqref{Wdefconstb2ii} (with $m$ replaced by $n$). Then commute $R_{s,\ell}(y)$ past $x^{n-k}$ using $f(y)x^{n-k}=x^{n-k}f(y-(n-k))$ and commute the remaining
$z^{r-k}$ past $y^m$ using $z^{r-k}y^m=(y-(r-k))^m z^{r-k}$. Grouping by the total loss index gives the recursion \eqref{Rrecconstb2ii}.

\item [\rm (5)] For $(x+y)^n$, write
$$
(x+y)^n = \sum_k x^k S_{n,k}(y)
$$ 
and multiply by $(x+y)$ on the left. Then
\begin{align*}
(x+y)^{n+1} = &\, x\sum_k x^k S_{n,k}(y)+y\sum_k x^k S_{n,k}(y) \\
= &\, \sum_k x^{k+1}S_{n,k}(y)+\sum_k x^k (y-k)S_{n,k}(y).
\end{align*}
using that $yx^k = x^k(y-k)$. Comparing coefficients yields the recursion for $S_{n,k}(y)$. 

For $(y+z)^n = \sum_k T_{n,k}(y)z^k$, multiply by $(y+z)$ on the left and use $z f(y)=f(y-1)z$ to obtain 
$$
T_{n+1,k}(y)=yT_{n,k}(y)+T_{n,k-1}(y-1).
$$
Finally, for $(x+z)^n$ write \eqref{Edefconstb2ii}, and expand
\[
(x+z)^{n+1}=(x+z)^n x+(x+z)^n z,
\]
and normal order $z^k x$ using $z^k x=\beta^k x z^k+b[k]_\beta z^{k-1}$. Collecting the coefficients of $x^i z^k$ gives \eqref{Erecconstb2ii}.
\end{enumerate}
\end{proof}

\end{document}
